\section{自诉案件的一审程序}
\begin{enumerate}
  \item 自诉案件永远别忘了公转自的情形！！
  \item 审查与受理。公安 \hz 检正在查，不予受理；其他基本同民诉，判断是不是符合条件。
  \item 当事人。
        \begin{enumerate}
          \item 可以调解、和解；可以撤诉、反诉。公转自不能调解、反诉。\al{和解、调解适用于自诉案件 \cuo}。
          \item 可分但不可反悔。
                \begin{enumerate}
                  \item 多名被告人可分。\al{乙丙丁，甲只告了乙}。
                  \item 多名自诉人可分。\al{甲乙丙，只有甲告了丁}。
                \end{enumerate}
        \end{enumerate}
  \item 法院。撤诉、错误立案、二审结案。
        \begin{enumerate}
          \item 自诉人缺乏罪证、证据不足错误立案。前有先后，后有或。前后处理不要混！
                \begin{enumerate}
                  \item 立案前发现，应当先说服撤诉 \yjt 拒不撤回的 \yjt 裁定不予受理。
                  \item 立案后发现，说服撤诉 \hz 直接裁定驳回起诉。\al{法院\xhx{应当说服}甲撤回起诉 \cuo}。
                \end{enumerate}
          \item 自诉人撤诉。
                \begin{enumerate}
                  \item 应当先审查自愿性 \yjt 再决定是否准许。\al{甲要求撤诉，应当裁定准许、裁定驳回起诉 \cuo}。
                  \item 部分自诉人撤诉或者被裁定按撤诉处理的，不影响案件的继续审理。
                  \item 传 2 次仍拒不到庭 \hz 中途退庭 \yjt 视为撤诉。同民诉，主动方不来撤诉处理。
                \end{enumerate}
          \item 被告下落不明。中止 \hz 决定逮捕。
          \item 受理后，通知公安移送材料，应当移送。
          \item \al{自诉人没有能力提供证据，可要求公安\xhx{提供协助} \dui{网络犯罪，取证技术要求较高}}。
          \item 二审结案方式。这里撤回自诉 = 撤回起诉！处理基本同民诉二审。
                \begin{enumerate}
                  \item 调解。必须制作调解书，送达后原判决视为撤销。
                  \item 和解。裁定撤回自诉时（地基），应当一并裁定撤销原一审判决（从头撤到尾）。
                  \item 提出反诉，应当告知另诉。民诉二审反诉可以调解。
                \end{enumerate}
          \item 审期。在押同公诉 2 + 1 + 3 + X。不在押审限同民诉。
        \end{enumerate}
\end{enumerate}
